\section{Conclusiones}

El proyecto cumplió su objetivo general de analizar las barreras de usabilidad, arquitectura de información y diseño de interacción presentes en WebSIS y de proponer una solución interactiva centrada en el estudiante. A través de un proceso de Diseño Centrado en el Usuario se logró pasar de un diagnóstico fundamentado en evidencia a un prototipo funcional evaluado con usuarios reales.

La investigación de usuarios, basada en la triangulación de encuesta, entrevistas y observación directa, demostró que las dificultades de WebSIS no son un problema individual de los estudiantes sino una consecuencia estructural del diseño de la interfaz. El sistema delega en el usuario una carga que debería asumir la propia plataforma, lo que se traduce en ansiedad, dependencia de ayuda externa y una relación predominantemente negativa con el sistema. La incorporación de la dimensión emocional al análisis permitió comprender que la experiencia de uso trasciende la eficiencia técnica.

La evaluación heurística del sistema actual permitió clasificar los problemas por severidad e identificar cuatro problemas catastróficos —ausencia de guía en la inscripción, mensajes de error incomprensibles, falta de compatibilidad móvil e inestabilidad técnica— como prioridad absoluta del rediseño. La arquitectura de la información se redefinió desde la lógica de tareas del estudiante y se validó con card sorting, confirmando su correspondencia con el modelo mental del usuario.

El prototipo funcional materializó la propuesta en rutas independientes que separan contextos de tarea, hacen visibles los estados críticos y emplean un lenguaje cercano al estudiante. Su evaluación, mediante pruebas de usabilidad y evaluación heurística, mostró que los flujos principales pueden completarse sin asistencia y detectó tres problemas de severidad menor a moderada. Estos problemas se corrigieron sobre el propio prototipo —prevención de choques de horario, mayor contraste del estado de matrícula y leyenda de las abreviaturas del Kardex—, completando un ciclo de detección, corrección y verificación propio del Diseño Centrado en el Usuario.

De este modo, el trabajo confirma la premisa central de la Interacción Humano-Computadora: un sistema no es funcional solo por operar correctamente, sino cuando el usuario logra alcanzar sus objetivos con efectividad, eficiencia y satisfacción. El rediseño propuesto, fundamentado en principios de usabilidad, leyes de la Gestalt y reducción de carga cognitiva, demuestra que es posible transformar una experiencia de uso frustrante en una más autónoma y confiable colocando al estudiante en el centro del proceso de diseño.

\subsection*{Limitaciones del estudio}

El proyecto presenta algunas limitaciones que conviene reconocer para enmarcar correctamente sus resultados:

\begin{itemize}
  \item \textbf{Tamaño de las muestras de evaluación.} La evaluación del prototipo se realizó con tres usuarios y el card sorting con ocho; aunque la literatura señala que muestras pequeñas detectan la mayoría de los problemas de usabilidad, una muestra mayor permitiría identificar problemas de menor frecuencia.
  \item \textbf{Alcance de carrera.} La investigación se concentró en estudiantes de Ingeniería Informática de la FCyT; otras carreras y facultades pueden presentar necesidades o contextos de uso distintos.
  \item \textbf{Naturaleza del prototipo.} La solución es un prototipo de alta fidelidad centrado en la interacción, no conectado a los sistemas reales de la UMSS; aspectos como rendimiento bajo carga, seguridad e integración con bases de datos institucionales quedan fuera del alcance.
  \item \textbf{Datos simulados.} La información mostrada (matrícula, Kardex, oferta de materias) es representativa pero simulada, por lo que no refleja la totalidad de casos reales posibles.
\end{itemize}

\subsection*{Trabajo futuro}

Una vez corregidos los problemas detectados en la evaluación, se plantea como continuación ampliar las pruebas de usabilidad a una muestra mayor y de carreras diversas, conectar el prototipo con datos reales para validar su comportamiento bajo carga, e incorporar de forma sistemática criterios de accesibilidad (navegación por teclado y compatibilidad con lectores de pantalla) para que la solución beneficie a la mayor cantidad posible de estudiantes.
